\section{A one-sided criterion for the symmetric powers of \texorpdfstring{$\Delta$}{Delta}}\label{sec:analytic}

Write $\Delta(z)=\sum_{n\ge1}\tau(n)e^{2\pi i nz}$ for Ramanujan's weight-twelve level-one newform. At each prime $p$, let
\[
 t_{p,m}=\alpha_p^m+\alpha_p^{-m},\qquad
 \alpha_p+\alpha_p^{-1}=\tau(p)p^{-11/2},\qquad |\alpha_p|=1 .
\]
The last equality is Deligne's bound. The level is one, so every finite Euler factor is unramified. Put $L_j(s)=L(s,\operatorname{Sym}^j\Delta)$ in unitary normalization, with $L_0=\zeta$. For $j\ge1$ symmetric-power automorphy and cuspidality are available for this non-CM form by Newton and Thorne \cite{nt}. We use the standard continuation, zero-free line $\Re s=1$, functional equation about $1/2$, and fixed-degree explicit formula for these functions \cite{nt,rtc}.

Let $P=\sum_{j=0}^{d}a_j C_j$ with $a_j>0$ for $0\le j\le d$, write $A=a_0$, and set
\begin{align*}
 S_P(x)&=\sum_{p^m\le x}(\log p)P(t_{p,m})
             \left(\frac1{p^m}-\frac1x\right),\\
 D_P(s)&=-\sum_{j=0}^{d}a_j\frac{L'_j}{L_j}(s),\qquad
 c_P=\lim_{s\to1}\left(D_P(s)-\frac{A}{s-1}\right)-A,\\
 R_P(x)&=S_P(x)-A\log x-c_P.
\end{align*}
Here $m$ ranges over positive integers. The identity
$C_j(t_{p,m})=\sum_{h=0}^j\alpha_p^{m(j-2h)}$
makes $D_P$ the Dirichlet series with coefficient $(\log p)P(t_{p,m})$ at $p^m$. One copy of the principal pole appears with weight $A$. The Haar extremizer supplies a $P$ which is also pointwise nonnegative, an additional property relevant to its arithmetic realization.

\begin{theorem}
For each fixed $P$ as above, the following are equivalent.
\begin{enumerate}
\item Every nontrivial zero of every $L_j$, $0\le j\le d$, lies on $\Re s=1/2$.
\item For every $0<\varepsilon<1/2$ there exist $K>0$ and $X\ge2$ such that
\[
 R_P(x)\le Kx^{-1/2+\varepsilon}\qquad(x\ge X).
\]
\end{enumerate}
The constants may depend on $P$ and $\varepsilon$.
\end{theorem}
\begin{proof}
For $\Re s>1$, termwise integration gives
\[
 \int_1^\infty R_P(x)x^{-s}\,dx
 =\frac{D_P(s)}{s(s-1)}-\frac{A}{(s-1)^2}-\frac{c_P}{s-1}.
\]
Our choice of $c_P$ removes both negative Laurent coefficients at $1$. At an off-line zero $\rho$ in the open strip, the residue of the right side is
\[
 \frac{\sum_j a_j\operatorname{ord}_{\rho}L_j}{\rho(1-\rho)}.
\]
It is nonzero at every such zero and strictly positive at a real zero. Common zeros contribute their orders additively.

Assume (2) and fix $\beta=1/2+\varepsilon$. For $x\ge X$, the function
$Kx^{\beta-1}-R_P(x)$ is nonnegative. Its Mellin integral with measure $x^{-s}\,dx$ converges for $\Re s>1$ by boundedness of $P$ and the elementary Chebyshev estimate. It has a meromorphic continuation in $\Re s>\beta$ by the displayed formula. If its real abscissa of convergence exceeded $\beta$, Landau's theorem for the transform of a nonnegative function would give a singularity there. The only possible singularities in this half-plane are simple poles from zeros of the $L_j$. At a real pole the residue of this transform is strictly negative, which would make the positive integral tend to $-\infty$ from the right. Removing the initial interval $[1,X]$ contributes an entire term. Thus the integral converges and is holomorphic on $\Re s>\beta$. Every zero there would produce a nonzero residue. Letting $\beta$ decrease to $1/2$ and applying the functional equations proves (1). This is the one-sided Mellin argument used in the repository \cite{source}; related transfers to weighted prime sums are discussed in the incoming reports \cite{deficit}.

Conversely, write $\psi_P(y)=\sum_{p^m\le y}\log p\,P(t_{p,m})$. Under (1), the fixed-degree explicit formula gives
$\psi_P(y)=Ay+O_P(\sqrt y\log^2(2y))$.
Indeed each zero term in the truncated formula is $O(\sqrt y/|\rho|)$, and the number of zeros up to height $T$ is $O_j(T\log(2T))$. Equivalently one may apply the fixed symmetric-power prime bound under GRH, then restore the higher prime powers by Chebyshev's estimate \cite{rtc}. Since $S_P(x)=\int_1^x\psi_P(y)y^{-2}\,dy$, the integral of $(\psi_P(y)-Ay)y^{-2}$ over $[1,\infty)$ converges and its tail is $O_P(x^{-1/2}\log^2(2x))$. The value of the integral is $c_P$, as follows by expanding $D_P(s)=s\int_1^\infty\psi_P(y)y^{-s-1}\,dy$ at $1$. This proves (2).
\end{proof}

\begin{proposition}
Let $R_P^{(1)}(x)$ use only $m=1$ in $S_P(x)$, and replace $c_P$ by
\[
 c_P^{(1)}=c_P-\sum_{p,\ m\ge2}\frac{\log p}{p^m}P(t_{p,m}).
\]
Then, unconditionally for this form,
\[
 R_P^{(1)}(x)-R_P(x)=\frac{2B_P+o(1)}{\sqrt{x}},
 \qquad B_P=\sum_{j=0}^{d}(-1)^j a_j.
\]
For the odd-degree extremizer $R_r$ of degree $2r-1$, this gives
$B_P=(r+1)(2r+1)/(6r)$. The one-sided criterion can therefore be stated using primes alone.
\end{proposition}
\begin{proof}
The correction to the centering constant is absolutely convergent. Put
$H(y)=\sum_{p^m\le y,\ m\ge2}\log p\,P(t_{p,m})$.
The terms $m\ge3$ are $O_P(y^{1/3}\log^2(2y))$ by the elementary Chebyshev estimate. For the squares, $t_{p,2}=t_{p,1}^2-2$. Sato--Tate for $\Delta$, the prime number theorem and partial summation give $H(y)=B_P\sqrt y+o(\sqrt y)$, with $B_P=\int P(t^2-2)\,d\mu(t)$. Indeed,
\[
 C_j(t^2-2)=\sum_{h=0}^j(-1)^h C_{2(j-h)}(t),
\]
by the character formula, so only the last term has nonzero Haar mean. Inserting the explicit coefficients of the extremizer and summing the two quadratic sequences yields the stated value. Partial summation gives
\[
 R_P^{(1)}(x)-R_P(x)=\frac{H(x)}{x}+\int_{(x,\infty)}\frac{dH(y)}{y}
   =\int_x^\infty\frac{H(y)}{y^2}\,dy,
\]
which proves the asymptotic, including its sign.
\end{proof}

For $d=1$, the pointwise weight is $P(t)=2+t$, and the criterion already tests the zeros of the degree-two $L$-function of $\Delta$ alongside those of $\zeta$. The missing arithmetic step is to prove the displayed upper bound for this trace-weighted remainder, at every sufficiently large $x$ for each $\varepsilon>0$. The repository's covolume comparison realizes a congruence and a quadratic-character weight. None of its short-form, spanning or determinant statements supplies the normalized trace $\tau(p)p^{-11/2}$. For higher $d$, the same question concerns the local weights $P(t_{p,m})$ constructed above.
